\section{问题三的模型建立与求解}

% ── 同 5_problem1 的结构规范 ──

\subsection{问题分析与模型选择}

% 为什么把本问建模为约束优化而非其他形式？
% 约束是硬约束还是软约束？目标函数的选择依据是什么？

将问题建模为约束优化问题：
\[
  \min \quad f(\mathbf{x}) = \sum_{i=1}^{n} c_i(x_i)
\]
\[
  \text{s.t.} \quad g_j(\mathbf{x}) \le 0, \quad h_k(\mathbf{x}) = 0
\]

选择该目标函数，而非[候选目标]，原因在于[题面依据或建模判断]。

\subsection{算法选择与实现}

采用遗传算法（种群规模 200，迭代 100 代）。选择遗传算法而非[候选算法]，
因为[选择理由：如离散决策空间/多峰问题/约束复杂]。

\subsection{求解结果}

得到最优策略方案：$x_1 = 0.75$，$x_2 = 1.20$，$x_3 = 0.45$，$x_4 = 0.88$。
所有约束均满足，验证方式：[说明如何验证可行性]。

\subsection{结果解释与机制}

% 必填：为什么最优解呈现这个结构？哪个约束是活跃的？
% 例：x₁ 的值受约束 g₁ 主导；移除该约束后目标值可降低 X%。

最优解的结构由[活跃约束]主导：[解释机制]。
当[某约束放松]时，目标值可改善[幅度]；说明[约束]是真正的瓶颈。

\subsection{适用边界}

本问结论在[参数范围/假设]下成立；若[关键假设变化]，则[结论如何变化]。

\subsection{直接答案}
